\section{Introduction}

Volatility is the main input to option pricing, margin setting, Value-at-Risk and position sizing,
but it cannot be observed directly. With intraday data it can be measured: the sum of squared 5-minute
returns gives a day's realized variance (RV), and forecasting RV has become the standard way to forecast
volatility \citep{andersen2003}. The benchmark is the HAR model of \citet{corsi2009}, a linear
regression of tomorrow's RV on its daily, weekly and monthly averages. It is simple and hard to beat.

Neural networks, in particular LSTMs, are often reported to beat HAR, but the gains are uneven and
the evidence comes mostly from liquid US stocks \citep{bucci2020,rahimikia2020,christensen2023}.
In less liquid stocks RV is a noisier measurement, which points two ways. A flexible model may find more
structure to exploit there (H1), or it may mostly fit noise and gain only in liquid stocks (H2).
Indian large caps span a wide range of liquidity --- average daily turnover in our sample runs from
about Rs~9 crore to Rs~1,360 crore --- so they are a natural place to test this.

\needspace{8\baselineskip}
We ask four questions:
\begin{enumerate}
\item Does an LSTM fed with HAR inputs (HAR-RV-LSTM) forecast next-day RV better than HAR-RV and log-HAR?
\item Does an LSTM fed with EGARCH features from daily returns do as well, and do EGARCH features add
anything to the HAR inputs?
\item Does the gain depend on the stock's liquidity (H1 against H2)?
\item Does it depend on market conditions, and is it stable over time and across training runs?
\end{enumerate}

\textbf{Main findings.} HAR-RV-LSTM is the most accurate model on every loss. Against classic HAR-RV it
reduces QLIKE by 15.4\%; against log-HAR, the fairer benchmark, by 1.6\%. Both differences are
significant, every training seed beats log-HAR, and the gain rises from 0.3\% in the most liquid stocks to
3.0\% in the least liquid, which supports H1. EGARCH-LSTM is worse than plain HAR-RV, and EGARCH features
add nothing to HAR-RV-LSTM. The gain is concentrated on calm and normal days.

\textbf{Contributions.} (i) A cleaned, reproducible 5-minute panel for 95 NSE stocks with documented
rules for pre-listing rows, bad prints and corporate actions. (ii) A liquidity-conditioned comparison
of econometric and deep-learning volatility models on an emerging market, with liquidity groups formed
only from training-period data. (iii) An honest account of how three implementation errors in our
first version --- uncleaned RV, a mis-scaled liquidity measure and EGARCH look-ahead --- had inflated
the apparent advantage of the neural models (Section~\ref{sec:versions}).
